\clearpage
\appendix
\thispagestyle{plain}
\vspace*{\fill}
\begin{center}
		\Huge\textbf{Appendix}
\end{center}
\vspace*{\fill}
\clearpage

\newpage

% \newgeometry{margin=0.75in}
\lstset{
    frame=none,
    aboveskip=2pt,
    belowskip=2pt,
    basicstyle=\ttfamily\footnotesize,
    breaklines=true
}

\chapter*{Appendix: Code Snippets}
\addcontentsline{toc}{chapter}{Appendix: Code Snippets}

\subsection{Alternating Least Squares Engine (als.py)}
\begin{lstlisting}[language=Python, caption={ALS Engine Implementation}]
import numpy as np
import time
from dataclasses import dataclass, field
from typing import List, Callable

# Convergence Tracker
@dataclass
class ALSLog:
  solver_name   : str
  train_rmse    : List[float] = field(default_factory=list)
  test_rmse     : List[float] = field(default_factory=list)
  train_loss    : List[float] = field(default_factory=list)
  residual_norm : List[float] = field(default_factory=list)
  cond_number   : List[float] = field(default_factory=list)
  elapsed_time  : List[float] = field(default_factory=list)

  def rmse_delta(self):
    tr = self.train_rmse
    return [abs(tr[i] - tr[i-1]) for i in range(1, len(tr))]

  def epoch_to_convergence(self, eps=1e-4):
    for i, d in enumerate(self.rmse_delta()):
      if d < eps: return i + 2
    return None

# Metric Helpers
def _rmse(R_true, U, V):
  mask = R_true > 0
  pred = (U @ V.T)[mask]
  return float(np.sqrt(np.mean((pred - R_true[mask]) ** 2)))

def _reg_loss(R, U, V, lam):
  mask = R > 0
  pred = (U @ V.T)[mask]
  return float(np.sum((pred - R[mask]) ** 2) 
               + lam * (np.sum(U ** 2) + np.sum(V ** 2)))

def _safe_cond(A):
  try:
    return float(np.linalg.cond(A))
  except np.linalg.LinAlgError:
    return float('inf')

# ALS Engine
def run_als(R_train, R_test, solver_fn, solver_name, k=10, lam=0.1, n_epochs=50, seed=42, verbose=True):
  log = ALSLog(solver_name=solver_name)
  np.random.seed(seed)
  n_users, n_items = R_train.shape
  lam_I = lam * np.eye(k)
  U = np.random.normal(0, 0.1, (n_users, k))
  V = np.random.normal(0, 0.1, (n_items, k))
  user_items = [np.where(R_train[u] > 0)[0] for u in range(n_users)]
  item_users = [np.where(R_train[:, i] > 0)[0] for i in range(n_items)]
  t0 = time.time()

  for epoch in range(n_epochs):
    ep_res, ep_cond = [], []
    for u in range(n_users):
      rated = user_items[u]
      if len(rated) == 0: continue
      V_Iu, r_u = V[rated], R_train[u, rated]
      A, b = V_Iu.T @ V_Iu + lam_I, V_Iu.T @ r_u
      x, res = solver_fn(A, b)
      U[u] = x
      ep_res.append(res); ep_cond.append(_safe_cond(A))

    for i in range(n_items):
      rated = item_users[i]
      if len(rated) == 0: continue
      U_Ii, r_i = U[rated], R_train[rated, i]
      A, b = U_Ii.T @ U_Ii + lam_I, U_Ii.T @ r_i
      x, res = solver_fn(A, b)
      V[i] = x
      ep_res.append(res); ep_cond.append(_safe_cond(A))

    log.train_rmse.append(_rmse(R_train, U, V))
    log.test_rmse.append(_rmse(R_test, U, V))
    log.train_loss.append(_reg_loss(R_train, U, V, lam))
    log.residual_norm.append(float(np.mean(ep_res)))
    log.cond_number.append(float(np.mean(ep_cond)))
    log.elapsed_time.append(time.time() - t0)
  return log
\end{lstlisting}

\subsection{Matrix Inverse Solver}
\begin{lstlisting}[language=Python, caption={Matrix Inverse Implementation}]
import numpy as np

def matrix_inverse(A, b):
    k = len(b)
    aug = np.zeros((k, 2 * k))
    for i in range(k):
        for j in range(k): aug[i, j] = float(A[i, j])
        aug[i, k + i] = 1.0

    for col in range(k):
        max_row = col + np.argmax(np.abs(aug[col:, col]))
        aug[[col, max_row]] = aug[[max_row, col]]
        pivot = aug[col, col]
        aug[col] /= pivot
        for row in range(k):
            if row == col: continue
            aug[row] -= aug[row, col] * aug[col]

    A_inv = aug[:, k:]
    x = A_inv @ b
    res = float(np.sqrt(np.sum((A @ x - b) ** 2)))
    return x, res
\end{lstlisting}

\subsection{Partial Pivoting Solver}
\begin{lstlisting}[language=Python, caption={Gaussian Elimination Implementation}]
import numpy as np

def partial_pivoting(A, b):
    k = len(b)
    M = np.column_stack((A.astype(float), b.astype(float)))
    for col in range(k):
        max_row = col + np.argmax(np.abs(M[col:, col]))
        M[[col, max_row]] = M[[max_row, col]]
        for row in range(col + 1, k):
            factor = M[row, col] / M[col, col]
            M[row, col:] -= factor * M[col, col:]

    x = np.zeros(k)
    for i in range(k - 1, -1, -1):
        x[i] = (M[i, k] - np.dot(M[i, i+1:k], x[i+1:k])) / M[i, i]
    res = float(np.sqrt(np.sum((A @ x - b) ** 2)))
    return x, res
\end{lstlisting}

\subsection{LU Decomposition Solver}
\begin{lstlisting}[language=Python, caption={LU Decomposition Implementation}]
import numpy as np

def lu_decomposition(A, b):
    k = len(b)
    L, U, P = np.eye(k), A.astype(float).copy(), np.eye(k)
    for col in range(k):
        max_row = col + np.argmax(np.abs(U[col:, col]))
        U[[col, max_row]], P[[col, max_row]] = U[[max_row, col]], P[[max_row, col]]
        L[[col, max_row], :col] = L[[max_row, col], :col]
        for row in range(col + 1, k):
            L[row, col] = U[row, col] / U[col, col]
            U[row, col:] -= L[row, col] * U[col, col:]

    pb = P @ b
    y = np.zeros(k)
    for i in range(k): y[i] = pb[i] - np.dot(L[i, :i], y[:i])
    x = np.zeros(k)
    for i in range(k - 1, -1, -1): x[i] = (y[i] - np.dot(U[i, i+1:], x[i+1:])) / U[i, i]
    res = float(np.sqrt(np.sum((A @ x - b) ** 2)))
    return x, res
\end{lstlisting}

\subsection{QR Decomposition Solver}
\begin{lstlisting}[language=Python, caption={QR Decomposition Implementation}]
import numpy as np

def qr_decomposition(A, b):
    k = len(b)
    R, b_hat = A.astype(float).copy(), b.astype(float).copy()
    for j in range(k):
        v = R[j:, j].copy()
        v[0] += (1.0 if v[0] >= 0.0 else -1.0) * np.sqrt(np.dot(v, v))
        norm_v = np.sqrt(np.dot(v, v))
        if norm_v == 0: continue
        v /= norm_v
        R[j:, j:] -= 2.0 * np.outer(v, np.dot(v, R[j:, j:]))
        b_hat[j:] -= 2.0 * v * np.dot(v, b_hat[j:])
    x = np.zeros(k)
    for i in range(k - 1, -1, -1): x[i] = (b_hat[i] - np.dot(R[i, i+1:], x[i+1:])) / R[i, i]
    res = float(np.sqrt(np.sum((A @ x - b) ** 2)))
    return x, res
\end{lstlisting}

\subsection{Cholesky Decomposition Solver}
\begin{lstlisting}[language=Python, caption={Cholesky Decomposition Implementation}]
import numpy as np

def cholesky_decomposition(A, b):
    k = len(b)
    L = np.zeros((k, k))
    for j in range(k):
        s = A[j, j] - np.dot(L[j, :j], L[j, :j])
        L[j, j] = np.sqrt(max(s, 0))
        for i in range(j + 1, k):
            L[i, j] = (A[i, j] - np.dot(L[i, :j], L[j, :j])) / L[j, j]
    y = np.zeros(k)
    for i in range(k): y[i] = (b[i] - np.dot(L[i, :i], y[:i])) / L[i, i]
    x = np.zeros(k)
    for i in range(k - 1, -1, -1): x[i] = (y[i] - np.dot(L[i+1:, i], x[i+1:])) / L[i, i]
    res = float(np.sqrt(np.sum((A @ x - b) ** 2)))
    return x, res
\end{lstlisting}

\restoregeometry
